% Generated draft: requires mathematical and source review.
\documentclass[11pt]{article}
\usepackage[margin=25mm]{geometry}
\usepackage{fontspec}
\setmainfont{Latin Modern Roman}
\usepackage{amsmath,amssymb,mathtools,unicode-math}
\setmathfont{Latin Modern Math}
\usepackage{enumitem,needspace,xurl,hyperref}
\setlength{\parindent}{0pt}
\setlength{\parskip}{5pt}
\setlength{\emergencystretch}{3em}
\newcommand{\N}{\mathbb{N}}
\newcommand{\Z}{\mathbb{Z}}
\newcommand{\Q}{\mathbb{Q}}
\newcommand{\R}{\mathbb{R}}
\newcommand{\C}{\mathbb{C}}
\newcommand{\F}{\mathbb{F}}
\newcommand{\A}{\mathbb{A}}
\newcommand{\PP}{\mathbb{P}}
\newcommand{\E}{\mathbb{E}}
\newcommand{\eps}{\varepsilon}
\newcommand{\dd}{\,\mathrm{d}}
\begin{document}
\section*{UnsolvedMath ID: 30005991}
\textbf{Source problem code:} catalogue entry
\subsection{Definitions and mathematical statement}
Let \(D\) denote the domain underlying a spectral optimal partition in the setting of the supplied source, and let \(\mathcal F_{\partial D}\subseteq\partial D\) denote its boundary free interface. Write \(N_{\partial}(x)\) for the limiting frequency at \(x\), whenever that limit is defined. Following the supplied formulation, set
\[
\operatorname{Reg}=\{x\in\mathcal F_{\partial D}:N_{\partial}(x)\text{ exists and }N_{\partial}(x)\in\Z\},
\qquad S=\mathcal F_{\partial D}\setminus\operatorname{Reg}.
\]
Here the integer-frequency criterion is retained exactly; it is not replaced by a frequency-two criterion.

The questions are whether, for every spectral optimal partition in the intended source setting,
\begin{enumerate}
\item \(\operatorname{Reg}\) is a \(C^1\) submanifold of codimension one in \(\partial D\);
\item \(S\) has codimension two in \(\partial D\).
\end{enumerate}
The supplied record does not specify the admissible domains, ambient dimension, number of partition components, spectral objective, admissibility conditions, boundary conditions, definition of the boundary free interface, or frequency formula. It also does not state hypotheses guaranteeing existence of the limiting frequency. Thus these questions retain their source-dependent setting and cannot be made into a fully specified universal assertion from this record alone. The second question retains the source wording: neither the notion of codimension nor whether it means exact codimension or a lower bound is specified. The first question likewise does not specify whether submanifolds with boundary are allowed.

The supplied dated assessment reports boundary regularity on the frequency-two regular set, while describing the full \(C^1\)-submanifold and codimension-two assertions as unresolved. It also reports a two-dimensional counterexample to an analogous codimension assertion for a broader nonregular set. The relationship between that broader set and \(S\) is not defined in the supplied record, so this report does not by itself identify a counterexample to the second question as formulated here. These are supplied literature qualifications, not independently verified mathematical conclusions.
\subsection{Short English statement}
For the boundary free interface of a spectral optimal partition, the supplied questions ask whether points with integer limiting frequency form a codimension-one \(C^1\) submanifold of the domain boundary, and whether the remaining interface has codimension two there. Essential hypotheses and definitions are missing. The supplied assessment labels the problem partially solved, citing frequency-two regularity and a separate qualification concerning a broader nonregular set.
\subsection{Sources}
\begin{itemize}
\item Supplied local UnsolvedMath registry record, permanent ID \(30005991\), title \emph{Regularity of Boundary Free Interfaces}, source problem code \texttt{catalogue entry}, in the collection \emph{Oberwolfach Reports Open Problems}. No individual proposer is supplied; the proposed year is recorded as 2024.
\item Supplied report citation: \emph{Calculus of Variations (2025)}, Oberwolfach Reports, DOI \texttt{10.4171/owr/2024/37}, pp.\ 2113--2178. Source URL and supplied literature-source link, labeled ``Oberwolfach source report'': \url{https://doi.org/10.4171/owr/2024/37}. The background describes these reports as documenting mathematical workshops and their open-problem discussions. No precise problem page, passage location, or individual speaker is supplied.
\item Literature assessment checked at \texttt{2026-08-22}: the supplied dated curation labels the problem \texttt{partially\_solved}, reports regularity on the frequency-two regular set, describes the full assertions as unresolved, and mentions a two-dimensional counterexample concerning a broader nonregular set. This is supplied dated curation, not an independently verified assessment; the record expressly requests independent verification of time-sensitive status and cited claims.
\item The registry labels the statement \texttt{corrected\_verified} and says it was reconstructed and checked against the official report. That verification claim is supplied metadata and has not been independently checked for this draft. The supplied report title contains 2025, while the DOI, proposed year, and year tag contain 2024; both are retained without resolving their bibliographic relationship.
\end{itemize}
\end{document}
